\input{../tex-inputs/latex-leseansicht-vorspann}

\section[Arthur, Olga Schnitzler und Emilie Sigmund an Gustav Schwarzkopf, 13. 9. 1904]{L04532 Arthur, Olga Schnitzler und Emilie Sigmund an Gustav Schwarzkopf, 13. 9. 1904}
\nopagebreak\mylabel{L04532v}
\rehead{ }\normalsize\beginnumbering\briefempfaengerindex{Schwarzkopf, Gustav@\textsc{Schwarzkopf, Gustav}!zzzSigmund, Emilie@\emph{von Emilie Sigmund}!1904-09-131@{13. 9. 1904}|(be}\briefempfaengerindex{Schwarzkopf, Gustav@\textsc{Schwarzkopf, Gustav}!zzzSchnitzler, Olga@\emph{von Olga Schnitzler}!1904-09-131@{13. 9. 1904}|(be}\briefempfaengerindex{Schwarzkopf, Gustav@\textsc{Schwarzkopf, Gustav}!zzzSchnitzler, Arthur@\emph{von Arthur Schnitzler}!1904-09-131@{13. 9. 1904}|(be}
\toendnotes[C]{\smallbreak\pagebreak[2]}
\correspDesc{Versand  durch Arthur Schnitzler, Olga Schnitzler, Emilie Sigmund am 13. 9. 1904 in St. Gilgen
\newline{}Übermittlung  am 14. 9. 1904 in St. Gilgen
\newline{}Zustellung  am 15. 9. 1904 in Wien
\newline{}Erhalt  durch Gustav Schwarzkopf am 15. 9. 1904 in Wien}\toendnotes[C]{\smallbreak}
\Standort{DLA, A:Schnitzler, HS.1985.1.1897.}
\physDesc{Bildpostkarte, 105 Zeichen
\newline{}Handschrift Arthur Schnitzler: Bleistift, deutsche Kurrent
\newline{}Handschrift Olga Schnitzler: Bleistift
\newline{}Handschrift Emilie Sigmund: Bleistift
\newline{}Versand: 1) Stempel: »\nobreak{}\oindex{St. Gilgen@\textbf{St. Gilgen}|pwk}St. Gilgen, {[}1{]}4. 9. 04\nobreak{}«.   2) Stempel: »\nobreak{}\oindex{I., Innere Stadt@\textbf{I., Innere Stadt}|pwk}{[}Wien{]} 1, 15. 9. 0\textcolor{gray}{4}, {[}Be{]}stellt\nobreak{}«. }\pstart{}{\pb}\textcolor{gray}{\textbf{An}}\pend{}\pstart{}\textsc{Herrn Gustav Schwarzkopf}\pend{}\pstart{}\textsc{Wien I}\oindex{Wien@\textbf{Wien}|pw}\pend{}\pstart{}\textsc{Tiefer Graben
                  23}\oindex{Wien@\textbf{Wien}!I., Innere Stadt@\textbf{I., Innere Stadt}!Tiefer Graben 23@\textbf{Tiefer Graben 23}|pw}.\pend{}{\bigskip}
\pstart
           \noindent{}\centering{}{\pb}\textcolor{gray}{\textbf{HOTEL {\kaufmannsund} PENSION LUEG ⋅ b.
                     St. Gilgen am Wolfgangsee}}\oindex{Hotel und Pension Lueg@\textbf{Hotel und Pension Lueg}|pw}\pend
           
\pstart
           \centering{}\textcolor{gray}{\textbf{DEPÔT DER STIEGL-BRAUEREI\orgindex{Stiegl-Brauerei@Stiegl-Brauerei|pw}
                     IN SALZBURG\oindex{Salzburg@\textbf{Salzburg}|pw}}}\pend
           
\pstart
           \centering{}\textcolor{gray}{\textbf{Bahn- {\kaufmannsund} Dampfschiff-Station}}\oindex{Bahn- und Dampfschiffstation St. Gilgen@\textbf{Bahn- und Dampfschiffstation St. Gilgen}|pw}\pend
           \vspace{1em}
\pstart
           \raggedleft{}{\pb}13. 9. 1904.\pend
           \vspace{0.5em}
\pstart
           Herzliche Grüße Ihnen und{\\[\baselineskip]}Dr. Max\pwindex{Schwarzkopf, Max 12.\,6.\,1857 Wien – 14.\,4.\,1928 ebd.@\textsc{Schwarzkopf, Max} (12.\,6.\,1857 Wien – 14.\,4.\,1928 ebd.)|pw}.\pend
           \leftskip=0em{}\selectlanguage{ngerman}\vspace{1em}\pstart {[}hs. O. Schnitzler:{]} \spacefill\mbox{OlgaS.}\pend{}\selectlanguage{ngerman}\vspace{1em}\pstart {[}hs. Sigmund:{]} \spacefill\mbox{E. Karin}\pend{}\selectlanguage{ngerman}\endnumbering\briefempfaengerindex{Schwarzkopf, Gustav@\textsc{Schwarzkopf, Gustav}!zzzSigmund, Emilie@\emph{von Emilie Sigmund}!1904-09-131@{13. 9. 1904}|)be}\briefempfaengerindex{Schwarzkopf, Gustav@\textsc{Schwarzkopf, Gustav}!zzzSchnitzler, Olga@\emph{von Olga Schnitzler}!1904-09-131@{13. 9. 1904}|)be}\briefempfaengerindex{Schwarzkopf, Gustav@\textsc{Schwarzkopf, Gustav}!zzzSchnitzler, Arthur@\emph{von Arthur Schnitzler}!1904-09-131@{13. 9. 1904}|)be}\mylabel{L04532h}
\begin{anhang}
\end{anhang}\newcommand{\dateiname}{L04532}\newcommand{\titel}{Arthur, Olga Schnitzler und Emilie Sigmund an Gustav Schwarzkopf, 13. 9. 1904}\newcommand{\editorInnen}{Herausgegeben von Jahnke, SelmaMüller, Martin Anton}\input{../tex-inputs/latex-leseansicht-abspann}
